\chapter{Fungsi Trigonometri}\label{ch:g12:trigo}

\omterm{def:g11:func:function}{Fungsi} $\cos$ dan $\sin$ mengubah geometri lingkaran menjadi
analisis. Didefinisikan dengan menggulung garis \omterm{def:g10:numbers:sets}{bilangan real} pada lingkaran
satuan, keduanya adalah purwarupa \omterm{def:g11:func:function}{fungsi} berkala, dan \omterm{def:g12:deriv:derivative}{turunannya} menjadikan
keduanya penyelesaian \omterm{def:g10:algebra:equation}{persamaan} osilasi $y'' = -y$.

\section{Lingkaran satuan}

\begin{definition}[Kosinus dan sinus]\label{def:g12:trigo:cossin}
Misalkan $\mathcal{C}$ lingkaran berjari-jari 1 yang berpusat di titik asal
suatu \omterm{def:g10:coordgeom:system}{sistem koordinat} \omterm{def:g10:coordgeom:system}{ortonormal}. Kepada setiap \omterm{def:g10:numbers:sets}{bilangan real} $t$ kaitkan
titik $M(t)$ yang diperoleh dengan menggulung panjang $\abs{t}$ sepanjang
$\mathcal{C}$ mulai dari titik $I(1,0)$, berlawanan arah jarum jam jika
$t \geq 0$ dan searah jarum jam jika sebaliknya. Maka \emph{$\cos t$ dan
$\sin t$ adalah koordinat $M(t)$}\index{kosinus}\index{sinus}:
\[
M(t) = (\cos t, \sin t).
\]
\end{definition}

\begin{omfigure}
\begin{tikzpicture}[scale=2.2]
  \draw[->] (-1.25,0) -- (1.35,0) node[right] {$x$};
  \draw[->] (0,-1.25) -- (0,1.35) node[above] {$y$};
  \draw (0,0) circle (1);
  \coordinate (O) at (0,0);
  \coordinate (M) at (50:1);
  \draw[thick,omDef] (O) -- (M);
  \fill (M) circle (0.6pt) node[above right] {$M(t)$};
  \fill (1,0) circle (0.6pt) node[below right] {$I$};
  \draw[dashed] (M) -- (M |- O) node[below,pos=1.05] {$\cos t$};
  \draw[dashed] (M) -- (M -| O) node[left,pos=1.02] {$\sin t$};
  \draw[->,omThm,thick] (0.35,0) arc (0:50:0.35);
  \node[omThm] at (25:0.5) {$t$};
\end{tikzpicture}
\end{omfigure}

Karena keliling lingkarannya $2\pi$, titik $M(t)$ tidak berubah ketika
$t$ bertambah sebesar $2\pi$; dan karena $M(t)$ terletak pada lingkaran satuan:

\begin{proposition}[Identitas dasar]\label{prop:g12:trigo:identities}
Untuk semua $t \in \R$ dan $k \in \Z$:
\[
\cos^2 t + \sin^2 t = 1, \qquad
\cos(t + 2k\pi) = \cos t, \qquad
\sin(t + 2k\pi) = \sin t,
\]
\[
\cos(-t) = \cos t, \qquad \sin(-t) = -\sin t,
\]
\[
\cos(\pi - t) = -\cos t, \quad \sin(\pi - t) = \sin t, \quad
\cos\left(\tfrac{\pi}{2} - t\right) = \sin t, \quad
\sin\left(\tfrac{\pi}{2} - t\right) = \cos t .
\]
\end{proposition}

\begin{proof}
Identitas pertama adalah \omterm{def:g10:algebra:equation}{persamaan} lingkaran satuan; sisanya menyatakan
kesetangkupan lingkaran itu: $M(-t)$ adalah pencerminan $M(t)$ pada
sumbu-$x$, $M(\pi - t)$ pada sumbu-$y$, dan $M(\frac\pi2 - t)$ pada garis
$y = x$.
\end{proof}

Nilai yang harus dihafal:
\begin{center}
\begin{tabular}{c|ccccc}
$t$ & $0$ & $\dfrac{\pi}{6}$ & $\dfrac{\pi}{4}$ & $\dfrac{\pi}{3}$ & $\dfrac{\pi}{2}$\\[6pt]
\hline\rule{0pt}{14pt}%
$\cos t$ & $1$ & $\dfrac{\sqrt3}{2}$ & $\dfrac{\sqrt2}{2}$ & $\dfrac12$ & $0$\\[6pt]
$\sin t$ & $0$ & $\dfrac12$ & $\dfrac{\sqrt2}{2}$ & $\dfrac{\sqrt3}{2}$ & $1$
\end{tabular}
\end{center}

\begin{proposition}[Rumus penjumlahan]\label{prop:g12:trigo:addition}
Untuk semua $a, b \in \R$:
\begin{align*}
\cos(a+b) &= \cos a \cos b - \sin a \sin b, &
\sin(a+b) &= \sin a \cos b + \cos a \sin b,\\
\cos(a-b) &= \cos a \cos b + \sin a \sin b, &
\sin(a-b) &= \sin a \cos b - \cos a \sin b.
\end{align*}
Khususnya $\cos 2a = 2\cos^2 a - 1 = 1 - 2\sin^2 a$ dan
$\sin 2a = 2 \sin a \cos a$.
\end{proposition}

\begin{proof}
\omterm{def:g11:vect:vector}{Vektor} $\vect{u} = (\cos a, \sin a)$ dan $\vect{v} = (\cos b, \sin b)$
adalah \omterm{def:g11:vect:vector}{vektor} satuan yang membentuk sudut $a - b$; \omterm{def:g11:scal:dot}{hasil kali skalarnya} sama
dengan $\cos a\cos b + \sin a \sin b$ (dari \omterm{def:g10:coordgeom:system}{koordinatnya}) sekaligus dengan
$\norm{\vect u}\,\norm{\vect v}\cos(a-b) = \cos(a - b)$ (dari geometrinya),
dan itulah rumus ketiga. Mengganti $b$ dengan $-b$ memberi rumus pertama;
rumus sinusnya menyusul dengan memakai $\sin x = \cos(\frac\pi2 - x)$. Rumus
pelipatduaannya adalah kasus $b = a$ digabung dengan $\cos^2 + \sin^2 = 1$.
\end{proof}

\section{Kajian analitis}

\begin{lemma}[Limit dasar]\label{lem:g12:trigo:sinxx}
\[
\lim_{t \to 0} \frac{\sin t}{t} = 1
\qquad\text{dan}\qquad
\lim_{t \to 0} \frac{\cos t - 1}{t} = 0 .
\]
\end{lemma}

\begin{proof}
Untuk $0 < t < \frac{\pi}{2}$, bandingkanlah tiga luas pada lingkaran satuan:
segitiga $OIM(t)$, juring lingkaran $OIM(t)$, dan segitiga siku-siku
beralas $OI$ dan bertinggi $\tan t$:
\[
\frac{\sin t}{2} \leq \frac{t}{2} \leq \frac{\tan t}{2}.
\]

\begin{omfigure}
\begin{tikzpicture}[scale=2.2]
  \draw[->] (-0.25,0) -- (1.5,0) node[right] {$x$};
  \draw[->] (0,-0.25) -- (0,1.45) node[above] {$y$};
  \fill[omDef!20] (0,0) -- (1,0) -- (50:1) -- cycle;
  \draw[gray, thin] (-10:1) arc (-10:80:1);
  \draw[omProp, thick] (0,0) -- (1,0) -- (1,1.1918) -- cycle;
  \draw[omDef] (0,0) -- (1,0) -- (50:1) -- cycle;
  \draw[omThm, thick] (1,0) arc (0:50:1);
  \draw[dashed] (50:1) -- (0.643,0);
  \node[left, font=\small] at (0.68,0.35) {$\sin t$};
  \node[omProp, right, font=\small] at (1,0.62) {$\tan t$};
  \node[omThm, font=\small] at (20:0.5) {$t$};
  \fill (50:1) circle (0.6pt);
  \node at (58:1.22) {$M(t)$};
  \fill (1,0) circle (0.6pt) node[below right] {$I$};
  \node[below left] at (0,0) {$O$};
\end{tikzpicture}

{\small Tiga luas yang bersarang: segitiga $OIM$ (biru, luasnya
$\frac{\sin t}{2}$), juring lingkarannya (dibatasi busur merah, luasnya
$\frac{t}{2}$), dan segitiga siku-siku bertinggi $\tan t$ (jingga, luasnya
$\frac{\tan t}{2}$).}
\end{omfigure}

Ketaksamaan pertama memberi $\frac{\sin t}{t} \leq 1$; yang kedua memberi
$\cos t \leq \frac{\sin t}{t}$. Ketika $t \to 0^+$, $\cos t \to 1$
(\omterm{def:g12:limcont:continuity}{kekontinuan}), dan teorema apit memberikan $\frac{\sin t}{t} \to 1$; sifat
ganjil-genapnya memperluas ini ke $t \to 0$. Untuk limit yang kedua,
\[
\frac{\cos t - 1}{t} = \frac{\cos^2 t - 1}{t(\cos t + 1)}
= -\frac{\sin t}{t}\cdot\frac{\sin t}{\cos t + 1}
\xrightarrow[t\to0]{} -1 \cdot \frac{0}{2} = 0. \qedhere
\]
\end{proof}

\begin{theorem}[Turunan sinus dan kosinus]\label{thm:g12:trigo:derivatives}
\omterm{def:g11:func:function}{Fungsi} $\sin$ dan $\cos$ \omterm{def:g12:deriv:derivative}{dapat diturunkan} pada $\R$, dengan
\[
(\sin)' = \cos, \qquad (\cos)' = -\sin .
\]
Secara lebih umum, $(\sin(ax+b))' = a\cos(ax+b)$ dan
$(\cos(ax+b))' = -a\sin(ax+b)$.
\end{theorem}

\begin{proof}
Menurut rumus penjumlahannya,
\[
\frac{\sin(x+h) - \sin x}{h}
= \sin x\,\frac{\cos h - 1}{h} + \cos x\,\frac{\sin h}{h}
\xrightarrow[h\to0]{} \sin x \cdot 0 + \cos x \cdot 1 = \cos x,
\]
dengan memakai \cref{lem:g12:trigo:sinxx}. Perhitungan untuk $\cos$ serupa,
dan bentuk umumnya menyusul dari aturan rantai.
\end{proof}

\begin{proposition}[Variasi]\label{prop:g12:trigo:variations}
Pada $\intcc{0}{\pi}$, $\cos$ turun dari $1$ ke $-1$. Pada
$\intcc{-\frac\pi2}{\frac\pi2}$, $\sin$ naik dari $-1$ ke $1$. Kedua \omterm{def:g11:func:function}{fungsi}
itu berkala $2\pi$ dan berdaerah hasil $\intcc{-1}{1}$; $\cos$ genap,
$\sin$ ganjil.
\end{proposition}

\begin{proof}
Pada $\intoo{0}{\pi}$, $(\cos)' = -\sin < 0$ sebab titik $M(t)$ berordinat
positif di sana; serupa pula $(\sin)' = \cos > 0$ pada
$\intoo{-\frac\pi2}{\frac\pi2}$. Sifat berkala dan sifat ganjil-genapnya
berasal dari \cref{prop:g12:trigo:identities}.
\end{proof}

\begin{omfigure}
\begin{tikzpicture}
\begin{axis}[
  width=12cm, height=4.5cm,
  axis lines=middle,
  xmin=-7, xmax=7, ymin=-1.4, ymax=1.4,
  xtick={-6.2832,-3.1416,3.1416,6.2832},
  xticklabels={$-2\pi$,$-\pi$,$\pi$,$2\pi$},
  ytick={-1,1},
  samples=200, domain=-7:7]
  \addplot[omDef,thick] {cos(deg(x))};
  \addplot[omThm,thick] {sin(deg(x))};
\end{axis}
\end{tikzpicture}

{\small Kurva $\cos$ (biru) dan $\sin$ (merah).}
\end{omfigure}

\begin{method}[Menyelesaikan $\cos x = a$ dan $\sin x = a$]\label{met:g12:trigo:equations}
Untuk $a \in \intcc{-1}{1}$, carilah satu penyelesaian khusus $\alpha$ (dari
tabel nilainya atau dengan kalkulator). Maka semua penyelesaiannya adalah:
\[
\cos x = \cos\alpha \iff x = \alpha + 2k\pi \text{ atau } x = -\alpha + 2k\pi,
\quad k \in \Z;
\]
\[
\sin x = \sin\alpha \iff x = \alpha + 2k\pi \text{ atau } x = \pi - \alpha + 2k\pi,
\quad k \in \Z.
\]
Untuk pertidaksamaannya, letakkanlah busur penyelesaiannya pada lingkaran
satuan lalu bacalah selangnya di dalam satu kala.
\end{method}

\begin{example}\label{ex:g12:trigo:equation}
Selesaikan $\cos x = \frac12$ di $\intoc{-\pi}{\pi}$: satu penyelesaian
khususnya $\frac\pi3$, jadi penyelesaiannya $x = \frac\pi3$ dan
$x = -\frac\pi3$. Di seluruh $\R$: $x = \pm\frac\pi3 + 2k\pi$, $k \in \Z$.
\end{example}

\section{Latihan}

\begin{exercise}[$\star$]\label{exo:g12:trigo:1}
Hitunglah $\cos\frac{2\pi}{3}$, $\sin\frac{5\pi}{6}$, $\cos\frac{7\pi}{4}$ dan
$\sin\left(-\frac{\pi}{3}\right)$ dengan memakai identitas pada
\cref{prop:g12:trigo:identities}.
\end{exercise}

\begin{exercise}[$\star$]\label{exo:g12:trigo:2}
Turunkanlah $f(x) = \sin^2 x$, $g(x) = \cos(3x) + x\sin x$ dan
$h(x) = \dfrac{\sin x}{2 + \cos x}$.
\end{exercise}

\begin{exercise}[$\star$]\label{exo:g12:trigo:3}
Selesaikan di $\intoc{-\pi}{\pi}$, lalu di $\R$:
\[
\text{(a) } \sin x = \frac{\sqrt3}{2};
\qquad
\text{(b) } \cos\left(2x\right) = \frac{\sqrt2}{2} .
\]
\end{exercise}

\begin{exercise}[$\star\star$]\label{exo:g12:trigo:4}
Dengan memakai rumus penjumlahannya, hitunglah nilai eksak
$\cos\frac{\pi}{12}$ dan $\sin\frac{\pi}{12}$.
(Petunjuk: $\frac{\pi}{12} = \frac{\pi}{3} - \frac{\pi}{4}$.)
\end{exercise}

\begin{exercise}[$\star\star$]\label{exo:g12:trigo:5}
Selesaikan di $\intcc{0}{2\pi}$ pertidaksamaan $2\sin^2 x - \sin x - 1 \geq 0$.
(Petunjuk: faktorkanlah bentuk kuadrat dalam $\sin x$ itu.)
\end{exercise}

\begin{exercise}[$\star\star$]\label{exo:g12:trigo:6}
Hitunglah limitnya
\[
\lim_{x\to0} \frac{\sin 3x}{x}, \qquad
\lim_{x\to0} \frac{1 - \cos x}{x^2}, \qquad
\lim_{x\to+\infty} x \sin\frac{1}{x}.
\]
\end{exercise}

\begin{exercise}[$\star\star$]\label{exo:g12:trigo:7}
Kajilah \omterm{def:g11:func:function}{fungsi} $f(x) = x + \cos x$ pada $\intcc{0}{2\pi}$: variasi dan
\omterm{def:g12:deriv:extremum}{nilai ekstremnya}. Tunjukkan bahwa $f$ naik pada $\R$ walaupun $f'$ bernilai
nol di tak berhingga banyak titik.
\end{exercise}

\begin{exercise}[$\star\star\star$]\label{exo:g12:trigo:8}
Misalkan $f(x) = A\cos x + B \sin x$ dengan $(A,B) \neq (0,0)$.
\begin{enumerate}
  \item Tunjukkan bahwa $f$ dapat ditulis $f(x) = R\cos(x - \varphi)$ dengan
        $R = \sqrt{A^2 + B^2}$ dan suatu $\varphi$ yang sesuai.
  \item Simpulkan nilai \omterm{def:g10:functions:extrema}{maksimum} dan \omterm{def:g10:functions:extrema}{minimum} $f$, lalu selesaikan
        $\cos x + \sin x = 1$ di $\intoc{-\pi}{\pi}$.
\end{enumerate}
\end{exercise}

\begin{exercise}[$\star\star\star$]\label{exo:g12:trigo:9}
Tunjukkan bahwa $\sin$ dan $\cos$ keduanya memenuhi \omterm{def:g10:algebra:equation}{persamaan} diferensial
$y'' = -y$. Sebaliknya, misalkan $y$ suatu penyelesaian $y'' = -y$ dengan
$y(0) = a$ dan $y'(0) = b$; tunjukkan bahwa $y = a\cos + b\sin$.
(Petunjuk: tinjaulah $g = (y - a\cos - b\sin)$ dan ``tenaga''
$E = g^2 + (g')^2$.)
\end{exercise}

\section{Soal: Rahasia bandul}

\begin{problem}[{Soal akhir pekan --- satu limit geometris menggerakkan
turunan sinus dan kosinus, dan dari keduanya seluruh
osilasi: pegas, bandul, layangan bunyi, dan meter yang nyaris
jadi}]
\label{pb:g12:trigo:1}
Setiap jam yang pernah berdetak, setiap dawai yang pernah berbunyi,
menuruti matematika yang sama: gaya pemulih yang sebanding dengan
simpangannya, sehingga \omterm{def:g10:algebra:equation}{persamaannya} $y'' = -\omega^2 y$, sehingga
\omterm{def:g12:trigo:cossin}{kosinus}. Di dasar semuanya duduk satu limit yang tampak tak berdosa,
$\frac{\sin x}{x} \to 1$ --- dijanjikan di kelas 9, dipakai oleh
\cref{thm:g12:trigo:derivatives}, dan dibuktikan di sini dengan gambar.
Soal ini mendaki dari limit itu sampai ke ayunan sebuah bandul
sepanjang satu meter, dan menjelaskan mengapa meter, tipis sekali
\omterm{def:g11:seq:arithmetic}{bedanya}, tidak didefinisikan olehnya.

\medskip
\noindent\textbf{Bagian I --- Limit dasarnya.}
\begin{enumerate}
  \item Pemanasan secara numeris: hitunglah $\frac{\sin x}{x}$ untuk
        $x = 0.1$ dan $x = 0.01$ (dalam \omterm{def:g11:trigo:radian}{radian}!).
  \item Apitannya, pada lingkaran satuan dengan
        $0 < x < \frac\pi2$: nyatakanlah dengan geometri dasar
        luas (i) segitiga yang bertitik sudut $O$, $A(1,0)$
        dan titik lingkaran $M(\cos x, \sin x)$; (ii) juring
        lingkaran $OAM$; (iii) segitiga siku-siku $OAT$
        dengan $T(1, \tan x)$.
  \item Dari apitan luas itu simpulkan
        $\sin x < x < \tan x$, lalu
        $\cos x < \frac{\sin x}{x} < 1$, dan simpulkan dengan
        teorema apit bahwa
        $\lim_{x \to 0} \frac{\sin x}{x} = 1$ (sifat ganjilnya mengurus
        $x < 0$).
  \item Simpulkan
        $\lim_{x \to 0} \frac{1 - \cos x}{x^2} = \frac12$
        (kalikan dengan sekawannya).
  \item Jelaskan mengapa limit ini adalah mesin
        \cref{thm:g12:trigo:derivatives} (berapakah
        $\sin'(0)$?), dan mengapa ia juga melunasi utang lama: dalam
        arti apa \emph{\omterm{def:g11:trigo:radian}{radian}} adalah satu-satunya satuan sudut yang
        membuat $\sin' = \cos$ berlaku tanpa tetapan tambahan?
\end{enumerate}

\medskip
\noindent\textbf{Bagian II --- Gerak harmonis.}
\begin{enumerate}[resume]
  \item Periksalah dengan aturan rantai bahwa
        $y(t) = R\cos(\omega t - \varphi)$ memenuhi
        $y'' = -\omega^2 y$.
  \item Sebuah massa pada pegas menuruti $y'' = -4y$, dengan
        $y(0) = 3$ dan $y'(0) = 8$. Dengan menerima (dari
        \cref{exo:g12:trigo:9}) bahwa semua penyelesaiannya berbentuk
        $A\cos 2t + B\sin 2t$, tentukanlah $A$ dan $B$.
  \item Tulislah ulang penyelesaiannya sebagai $R\cos(2t - \varphi)$
        (\cref{exo:g12:trigo:8}): berikan amplitudo $R$, kala
        ayunnya, dan fase $\varphi$ (tiga angka di belakang titik).
  \item Pada saat manakah massa itu untuk pertama kalinya melewati
        kedudukan seimbang $y = 0$?
  \item Bandulnya: sebuah bandul pada tali sepanjang $L$ menuruti
        \emph{secara eksak} $\theta'' = -\frac gL \sin\theta$ ---
        tak terselesaikan dengan \omterm{def:g11:func:function}{fungsi} dasar. Untuk ayunan kecil,
        gantilah $\sin\theta$ dengan $\theta$ (Bagian I membenarkannya!)
        lalu simpulkan kala ayunnya
        $T = 2\pi\sqrt{\frac Lg}$. Hitunglah $T$ untuk
        $L = 1$ m ($g = 9.81$): penemuan masyhur Galileo ---
        pada apakah rumus itu \emph{tidak} bergantung?
  \item Bandul sekon: panjang berapakah yang memberi
        $T = 2$ sekon tepat? Akademi Ilmu Pengetahuan Prancis pada tahun
        1791 ragu antara panjang itu dan sepersepuluh juta
        seperempat meridian untuk mendefinisikan \emph{meter} --- berapa
        milimeterkah \omterm{def:g11:seq:arithmetic}{beda} kedua calon itu?
\end{enumerate}

\medskip
\noindent\textbf{Bagian III --- Rumus penjumlahan beraksi.}
\begin{enumerate}[resume]
  \item Tunjukkan dengan rumus penjumlahannya
        (\cref{prop:g12:trigo:addition}) bahwa
        $\cos\left(x + \frac\pi2\right) = -\sin x$ dan
        $\sin\left(x + \frac\pi2\right) = \cos x$: penurunan
        secara melingkar adalah seperempat putaran --- periksalah
        keselarasannya dengan \cref{thm:g12:trigo:derivatives}.
  \item Turunkanlah identitas hasil-kali-ke-jumlah
        $\cos a \cos b = \frac12\left(\cos(a - b) +
        \cos(a + b)\right)$ dan pelinearan
        $\cos^2 x = \frac{1 + \cos 2x}{2}$ (bab pengintegralan
        akan berterima kasih).
  \item Layangan bunyi: turunkanlah
        $\sin p + \sin q = 2 \sin\frac{p + q}{2}
        \cos\frac{p - q}{2}$, lalu terapkan pada dua garpu tala
        pada $440$ dan $444$ Hz: nada apakah yang didengar telinga, dan
        berapa denyut tiap sekonnya? (Beginilah orkestra
        menyetel nadanya.)
  \item Turunkanlah identitas sudut rangkap tiga
        $\cos 3x = 4\cos^3 x - 3\cos x$
        (tulislah $3x = 2x + x$). Dengan mengambil $x = 20^\circ$,
        \omterm{def:g10:algebra:equation}{persamaan} pangkat tiga apakah yang harus dipenuhi
        $\cos 20^\circ$? (Bahwa \omterm{def:g10:algebra:equation}{persamaan} pangkat tiga itu tak dapat
        diselesaikan dengan akar kuadrat saja justru itulah sebabnya
        pembagian sudut menjadi tiga mengalahkan penggaris dan jangka
        selama dua ribu tahun --- kisah lengkapnya
        diceritakan dalam jilid universitas.)
  \item Sebuah arus bolak-balik terbaca
        $I(t) = 3\cos(100\pi t) + 3\sqrt3 \sin(100\pi t)$.
        Tulislah ulang sebagai $R\cos(100\pi t - \varphi)$: berapakah
        arus puncaknya dan fasenya?
\end{enumerate}

\medskip
\noindent\textbf{Bagian IV --- Dunia yang berayun.}
\begin{enumerate}[resume]
  \item Kala ayun: berapakah kala (terkecil) dari
        $\sin^2 t$ (pakailah pertanyaan 13)? Dan dari
        $\sin(2t) + \sin(3t)$?
  \item Bandul yang sesungguhnya teredam:
        $y(t) = \eu^{-t/10}\cos(2\pi t)$. Lukiskan geraknya,
        hitunglah nilai selubungnya di $t = 10$ (seorang kawan lama
        muncul kembali, \cref{pb:g12:exp:1}), lalu sebutkan bab
        tempat \omterm{def:g10:algebra:equation}{persamaan} dengan penyelesaian seperti itu diselesaikan.
  \item Mengapa \omterm{def:g12:trigo:cossin}{sinus} dan \omterm{def:g12:trigo:cossin}{kosinus} menguasai setiap getaran? Berikan
        kamus fisika-ke-matematikanya dalam dua kalimat
        (gaya pemulih yang sebanding dengan simpangannya
        $\rightarrow$ \omterm{def:g10:algebra:equation}{persamaan} yang mana $\rightarrow$ penyelesaian
        yang mana), lalu katakan apa sumbangan \cref{exo:g12:trigo:9}
        bagi pernyataan itu.
  \item Penutup: busur bab ini dalam empat langkah --- satu
        limit geometris, dua \omterm{def:g12:deriv:derivative}{turunan}, satu \omterm{def:g10:algebra:equation}{persamaan}
        diferensial, seluruh osilasi. Dan tanda bintang yang jujur:
        untuk ayunan besar, kala bandul yang sebenarnya melibatkan
        integral eliptis, yang dihitung secepat kilat oleh\,\dots
        rata-rata aritmetika--geometri Gauss, hidangan penutup
        \cref{pb:g12:seq:1}. Benang-benang seri ini menyimpul dirinya.

\end{enumerate}
\end{problem}
